% ===========================================================================
% Dataset section -- regenerated 2026-09-14 from the deduplicated corpus.
% Placeholders marked \TODO{} are values the dynamic typing campaign has not
% produced yet.  Every other number is measured; see docs/paper/NUMBERS.md for
% provenance.
% ===========================================================================

We began with a benign sample dataset to establish our baseline and conduct a
preliminary study. This dataset included samples from WinGet (Windows Package
Manager)~\cite{winget} and was packed with \textbf{41} off-the-shelf packer
families comprising \textbf{107} variants, all targeting 32-bit Intel i386 PE32
executables; see \Cref{tab:packer_families}.

\begin{table}[t]
\centering
\caption{Packer families and variant counts in the corpus. The corpus
comprises \textbf{107} packer variants across \textbf{41} distinct
families.}
\label{tab:packer_families}
\footnotesize
\setlength{\tabcolsep}{4pt}
\begin{tabular}{@{}lr@{\hspace{1.6em}}lr@{}}
\toprule
Family & \#V & Family & \#V \\
\midrule
UPX                 & 50 & hxor\_packer     & 1 \\
UPX\_Scrambler      & 6  & jdpack           & 1 \\
yoda\_protector     & 5  & mew              & 1 \\
amber               & 3  & molebox          & 1 \\
petite              & 3  & npack            & 1 \\
hyperion            & 2  & nspack           & 1 \\
kkrunchy            & 2  & obsidium         & 1 \\
mpress              & 2  & packman          & 1 \\
yoda\_crypter       & 2  & pe\_diminisher   & 1 \\
acprotect           & 1  & pelock           & 1 \\
alienyze\_protector & 1  & pepacker         & 1 \\
alushpacker         & 1  & pezor            & 1 \\
armadillo           & 1  & rlpack           & 1 \\
asm\_guard          & 1  & shrinker         & 1 \\
astral\_pe          & 1  & telock           & 1 \\
beroexepacker       & 1  & hackupx          & 1 \\
enigma              & 1  & upack            & 1 \\
eronona             & 1  & zprotect         & 1 \\
exe32pack           & 1  & xcomp            & 1 \\
fsg                 & 1  &                  &   \\
\midrule
\multicolumn{4}{@{}l}{\textbf{Total: 41 families, 107 variants}} \\
\bottomrule
\end{tabular}
\end{table}

For replication of real-world packed malware for studying purposes, we curated a
dataset featuring the popular malware tool VirusTotal and its Folder of Malware
Samples for Academic Researchers~\cite{virustotal}, low-entropy malware
from~\cite{mantovani2020prevalence}, and DikeDataset~\cite{dikedataset}. After
removing \textbf{610} exact-duplicate binaries identified by SHA-256 over the
extracted samples, the dataset totals \textbf{44,740} distinct samples. Then, to
distinguish packed samples from unpacked malware samples, we verified packedness
using packer identification software, including DIE, PEID, CAPA, YARA, and
Shannon Entropy Score~\cite{yara,die,manalyze,capa,pefile,lyda2007entropy,peid}.
Although the software may have disagreed on the packer type, an executable was
added to our dataset if all the tools recognized the sample as packed. For the
residual samples that these faster detectors did not flag, we additionally
performed dynamic analysis based on the run-time unpacking
behaviors~\cite{ugartepedrero2015sok}, tracing written-then-executed memory
regions; this resulted in a total \textbf{72.9\%} of samples being flagged as
packed.

We determined each sample's format by parsing PE and ELF headers statically,
without execution: PE32 with an \texttt{i386} machine field is Win32, PE32+ with
x86-64 is Win64, and any PE carrying a CLR data directory is managed (.NET).
\Cref{tab:filetypes} gives the distribution. \textbf{This study is scoped to the
42,656 Win32 (PE32/i386) samples, of which 31,211 are packed.} Managed assemblies
do not unpack through the native write-then-execute mechanism our oracle observes,
and the remaining formats cannot execute in the i386 guest.

\begin{table}[t]
\centering
\caption{File-format distribution of the deduplicated malware corpus
(\textbf{44,740} distinct samples), determined by static header parsing.
``Packed'' is the subset flagged by the detector ensemble described above.}
\label{tab:filetypes}
\footnotesize
\begin{tabular}{@{}lrrr@{}}
\toprule
Format & Samples & Share & Packed \\
\midrule
Win32 (PE32, i386)      & 42,656 & 95.3\% & 31,211 \\
Win64 (PE32+, x86-64)   &  1,114 &  2.5\% &    761 \\
.NET (CLR)              &    909 &  2.0\% &    735 \\
MS-DOS (MZ, no PE)      &     55 &  0.1\% &      3 \\
PE32 (non-i386, ARM)    &      4 &  0.0\% &      3 \\
Unclassified            &      2 &  0.0\% &      1 \\
\midrule
\textbf{Total}          & \textbf{44,740} & \textbf{100\%} & \textbf{32,714} \\
\bottomrule
\end{tabular}
\end{table}

\Cref{tab:typedist} reports the Ugarte Type distribution. Every benign condition
is typed from its own dynamic trace. A malware sample inherits a Type only when
its static packer label exactly matches a corpus family typed dynamically;
otherwise it requires its own trace. Of \textbf{8,224} Win32 packed samples typed
so far, \textbf{830} are dynamic and \textbf{7,394} inherited. The campaign over
the remaining \textbf{22,987} is ongoing; those entries are blank.

\begin{table}[t]
\centering
\caption{Ugarte Type distribution. Benign counts are packer \emph{conditions},
each typed from its own dynamic trace. Malware counts are Win32 packed
\emph{samples}. Blank entries await the ongoing dynamic typing campaign.}
\label{tab:typedist}
\footnotesize
\begin{tabular}{@{}lrrrr@{}}
\toprule
 & \multicolumn{1}{c}{Benign} & \multicolumn{3}{c}{Malware (Win32 packed)} \\
\cmidrule(lr){2-2}\cmidrule(lr){3-5}
Type & Conditions & Dynamic & Inherited & Total \\
\midrule
TYPE\_0     &   1 & \TODO{} & \TODO{} & \TODO{} \\
TYPE\_I     &  76 & \TODO{} & \TODO{} &  1,634 \\
TYPE\_II    &   7 & \TODO{} & \TODO{} &    129 \\
TYPE\_III   &   5 & \TODO{} & \TODO{} &    218 \\
TYPE\_IV    &  11 & \TODO{} & \TODO{} &  5,840 \\
TYPE\_V-F   &   2 & \TODO{} & \TODO{} &    121 \\
TYPE\_VI-F  &   5 & \TODO{} & \TODO{} &    282 \\
\midrule
Typed       & \textbf{112} & \textbf{830} & \textbf{7,394} & \textbf{8,224} \\
Untyped     &   0 & \multicolumn{3}{r}{\textbf{22,987}} \\
\midrule
\textbf{Total} & \textbf{112} & \multicolumn{3}{r}{\textbf{31,211}} \\
\bottomrule
\end{tabular}
\end{table}

\subsection{A held-out test set from EMBER2024}

We adopt the Win32 portion of EMBER2024~\cite{ember2024} as a held-out test set.
It publishes features and labels, but not binaries, for files first submitted to
VirusTotal between 2023-09-24 and 2024-12-14, with packer labels assigned by
ClarAVy~\cite{claravy}. Two properties make it suitable: an independently produced
packer label to evaluate against, and disjointness from our corpus ---
cross-referencing all \textbf{44,738} corpus SHA-256 digests against the
\textbf{1,919,994} EMBER Win32 digests yields \textbf{zero} intersection.

Two defects in the published release affect any reuse. First, the Win32 archives
contain every record twice --- 3,120,000 train and 720,000 test lines for
1,560,000 and 359,994 distinct digests, differing only in \texttt{caps},
\texttt{ttps}, and \texttt{mbc} --- so the reference loader double-counts; we
deduplicate by SHA-256, recovering EMBER's documented totals. Second, the packer
distribution is not stationary across the split: \texttt{themida} rises from
11.8\% to 34.6\% of labels while \texttt{upx} falls from 47.0\% to 29.7\%. We
therefore report per-family rather than aggregate results.

Of the 1,919,994 EMBER Win32 samples, \textbf{71,173} (3.7\%) carry a packer
label, giving \textbf{71,779} occurrences across \textbf{49} packer names.
\textbf{18} names match a corpus family typed dynamically, covering
\textbf{46,336} occurrences (\textbf{64.6\%}); the largest unmatched are
\texttt{vmprotect} (10,582) and \texttt{nsis} (10,344), together 29.1\%. We apply
the same dynamic analysis to the \TODO{} EMBER samples obtainable as binaries.
\Cref{tab:ember} reports both distributions.

\begin{table}[t]
\centering
\caption{EMBER2024 Win32 held-out test set. ``Inherited'' counts label
occurrences whose packer name matches a corpus family with a dynamically
established Type. ``Dynamic'' counts samples we obtained as binaries and typed
with the same oracle; blank entries await that campaign.}
\label{tab:ember}
\footnotesize
\begin{tabular}{@{}lrr@{}}
\toprule
Type & Inherited occurrences & Dynamically typed \\
\midrule
TYPE\_I     & 43,132 & \TODO{} \\
TYPE\_II    &    601 & \TODO{} \\
TYPE\_III   &  1,382 & \TODO{} \\
TYPE\_IV    &    827 & \TODO{} \\
TYPE\_V-F   &      7 & \TODO{} \\
TYPE\_VI-F  &    387 & \TODO{} \\
\midrule
Matched to a corpus family & \textbf{46,336} & \TODO{} \\
Unmatched packer names     & \textbf{25,443} & --- \\
\midrule
\textbf{Total labelled}    & \textbf{71,779} & \TODO{} \\
\bottomrule
\end{tabular}
\end{table}
